\chapter{进化不只有自然选择}

\begin{kaichang}
找一个不透明的罐子，装进去 50 颗红弹珠和 50 颗黑弹珠，摇匀。现在闭上眼睛，抓出一把——就 10 颗。数数看，是 5 红 5 黑吗？

多半不是。可能是 6 红 4 黑，也可能是 3 红 7 黑。你没有任何偏好，手指也不认识颜色，可抓出来的比例就是和罐子里不一样。再抓一把放回去之前，先记下这次的“成绩”。如果这个游戏一代一代玩下去——每一代都只有 10 颗弹珠能“活下来”并生出下一代的 10 颗——几百年、几百万代之后，罐子里会变成什么样？

先别急着说“最后还是一半一半，公平嘛”。这一章结束时你会发现，达尔文之外，还站着一位从不说话、却威力巨大的“玩家”：运气。自然选择确实重要（第 76 章已经讲过），但进化的牌桌上，发牌的常常是偶然。
\end{kaichang}

\section{先想清楚：什么才算“进化”}

上一章我们给进化下过一个很朴素的定义：\textbf{群体中等位基因的频率，随世代发生改变}。注意这个定义里没有一个字提到“选择”。它只说频率变了，没说为什么变。

这就好像水位。湖水上涨，可能是上游下了暴雨（相当于选择：环境施加方向性的压力），但也可能只是有人开闸放水、有人抽水灌溉、甚至只是测量那天起了风浪。水位变化是事实，原因却可以有很多种。等位基因频率也一样——任何能改变它的机制，都是进化的发动机。自然选择只是其中马力最大、也最容易被看见的一台。

生物学家数得出来的发动机至少有这么几台：\textbf{突变}不断制造新的变异；\textbf{遗传漂变}让频率纯粹因为抽样运气而起伏；\textbf{基因流}在不同群体之间搬运等位基因；\textbf{性选择}专门改写“谁能留下后代”的规则。它们常常同时运转，有的把群体推向适应，有的纯粹制造噪音，还有的专门给噪音“记日记”。这一章，我们一台一台拆开看。

先打个预防针：这些机制不是达尔文理论的“反对派”。它们全都建立在同一座地基上——第 71 章的遗传规律、第 72 章的突变、第 76 章的群体视角。达尔文画出了进化大厦的主梁，而这一章讲的是让大厦真正立得住的其余骨架。

\section{突变：一切变异的源头活水}

没有原材料，什么选择、漂变都无从谈起。所有等位基因，追到底都是某一次突变的后代。

第 72 章已经拆过突变的种类：碱基替换、插入、缺失、重复、倒位、易位。这里只强调两个常被忽略的要点。

第一，突变的方向是\textbf{随机的}——不是“环境需要什么就突变什么”。细菌不会因为培养皿里有抗生素，就“定向”制造耐药突变。耐药突变一直在以极低的概率随机出现，抗生素只是杀死没有它的个体，让少数幸运儿留下来（第 76 章的桦尺蛾和耐药菌例子讲的就是这个筛选过程）。突变负责出题，选择负责阅卷，出题人并不知道标准答案。

第二，突变的速度虽然慢，架不住基数大。来做一个真实的数量级估算：人类基因组约有 $3\times 10^{9}$ 个碱基对，每个碱基每一代的突变率大约是 $10^{-8}$ 量级。两项相乘：

\[
3\times 10^{9}\ \times\ 10^{-8}\ \approx\ 30
\]

也就是说，平均每个新生儿都带着大约几十处父母身上没有的新突变（更精细的测量给出的数字在 60 上下，因为还要计入父母双方各自的贡献）。听起来很少？把这个数乘以全世界每年出生的约 1.3 亿婴儿：人类这个物种\textbf{每一代}都会产生几十亿处全新的突变。基因组上几乎每一个可以出错的位置，每一代都有人在出错。自然选择从来不必“等”变异出现——货架上永远摆得满满当当。

绝大多数突变落进两个筐：要么中性（不改变生存和繁殖的机会，比如发生在不编码蛋白质、也不影响调控的序列里），要么有害（破坏了正常功能）。真正“有利”的是极少数，而且离开了具体环境谈不上有利——让镰刀型贫血的等位基因在疟疾流行区是保命符，换个地方就只是病（第 76 章）。这正是第 72 章那句结论的进化版：突变不是错误，也不全是创新，它是材料。

\section{漂变：运气也是进化力量}

回到开场的弹珠罐。把红弹珠想成某个等位基因 $A$，黑弹珠想成 $a$，罐子里 50 对 50，频率各一半。抓出 10 颗当“下一代的亲本”，其实就是一次随机抽样。抽出来的比例很少恰好是一半——下一代的频率就这样莫名其妙地变了。再抽一代，又变一次。这个过程没有任何“适者生存”的成分，纯粹是概率在起伏，却完全符合进化的定义：频率变了。生物学家给它起的名字叫\textbf{遗传漂变}。

你可能觉得这听起来像小孩子的游戏，不值一提。那先看一个数学直觉。抛一枚均匀硬币 10 次，得到 7 次正面毫不稀奇，偏差 20\% 稀松平常；抛 10000 次，正面比例几乎一定死死钉在 50\% 附近，偏差很难超过 1\%。样本越小，运气说话越大声；样本越大，运气的声音被平均掉。这就是大数定律的口头版。

把“抛硬币次数”换成“群体里的个体数”，结论直接搬过来：\textbf{遗传漂变在小群体中威力巨大，在大群体中几乎听不见。}一个只有几十只个体的孤岛鸟群，某个等位基因可能因为几只好运的鸟多生了几个蛋，就在几代之内从稀有变成普遍，甚至彻底固定下来（频率变成 100\%）；也可能反过来，无声无息地消失。而这一切和“它有没有用”毫无关系。

\begin{qiaoliang}
抽样波动到底有多大？如果群体里等位基因 $A$ 的频率是 $p$，从中随机抽取 $N$ 个个体繁殖下一代，下一代频率的标准差大约是
\[
\sqrt{\frac{p(1-p)}{N}}
\]
代入数字感受一下：$p=0.5$ 时，$N=10$ 的小群体，标准差约 $0.16$——下一代频率漂到 $0.34$ 或 $0.66$ 都是家常便饭；$N=10^{6}$ 的大群体，标准差只有约 $0.0005$，漂到小数点后三位都费劲。同一个随机过程，在不同规模的群体里完全是两种脾气。这就是为什么保护生物学张口闭口都是“群体数量”：个体越少，基因库就越像那罐只抓一把的弹珠。
\end{qiaoliang}

\begin{shiyan}
\textbf{豆子漂变模拟}（在家可做，约 20 分钟）

材料：两种颜色的干豆子（红豆、绿豆各 20 粒）、一个不透明的袋子或小盒子、纸和笔。

步骤：
\begin{enumerate}[itemsep=1pt]
\item 把 20 红 20 绿放进袋子，摇匀。这是“第 0 代”，红豆等位基因频率 50\%。
\item 不看袋子，摸出 10 粒。这 10 粒代表“有幸繁殖下一代的个体”。
\item 数出红豆的粒数，记为下一代的红豆频率（例如 6 粒就是 60\%）。然后按这个新比例“重建”一个 20 粒的群体（12 红 8 绿），放回袋子。摸出的豆子记录完即可放回。
\item 重复第 2、3 步，做 10 代以上，把每代的红豆频率连成一条折线。
\item 进阶：改用摸 20 粒重建 40 粒群体（相当于群体变大），再跑一遍，比较两条折线。
\end{enumerate}

观察点：小群体那条折线是不是上下乱跳，甚至某一代红豆彻底消失或占到全部？大群体那条是不是老实得多？同一种颜色一旦占到 100\%，还能再变回来吗？（不能——没有突变补充的话，丢失的等位基因就永远丢失了。）

安全提示：干豆子不是零食，请勿食用；结束后收好，避免幼儿和宠物误吞。
\end{shiyan}

\section{奠基者与瓶颈：漂变的两个著名场景}

自然界的“小群体”不是抽象数字，它们常常以两种戏剧性方式出现。

\textbf{奠基者效应}：一小撮个体离开大群体，去开辟新天地。这个新群体的基因库，就是出发时那“一把弹珠”的拷贝——不管原来群体多大多多样，新群体只带着这几位的等位基因。真实案例在西太平洋的平格拉普环礁：1775 年一场台风加饥荒过后，岛上只剩约 20 名幸存者，其中一位恰好携带一种导致全色盲（完全看不见颜色）的隐性等位基因。在普通人群里，这个等位基因的频率低到万分之一量级；而今天平格拉普岛上大约每十个人就有一人患全色盲。不是这个基因“适合”海岛生活，它只是搭上了幸存者的船。类似的还有北美阿米什人社区——几百年前少数欧洲移民的后代，某些罕见遗传病的频率远高于一般人群，根源同样是奠基者人数太少。

\textbf{瓶颈效应}：群体没有搬家，而是被灾难狠狠地“筛”了一遍。幸存者的等位基因构成，从此成为整个物种的全部家底。北象海豹在 19 世纪被猎杀到只剩几十头，如今数量恢复到十几万，却个个都像近亲：遗传学家测它们的蛋白质变异，几乎测不到个体差异。中国的朱鹮是更贴近的例子：1981 年在陕西洋县找到全世界仅存的 7 只野生朱鹮，今天种群恢复到数千只，但它们全是那 7 只的后代。数量回来了，基因多样性却没有——这直接引出保护生物学的一条核心原则：\textbf{保护一个物种，要保护的不只是个体数量，还有它的遗传多样性。}数量可以繁殖回来，丢失的等位基因只能靠漫长得多的突变慢慢重新积累。

瓶颈留下的隐患很现实：多样性低的群体像一支全队同一个打法的球队，遇到没见过的新疾病或环境剧变，可能全线崩溃。今天全世界的朱鹮保育工作里，谱系登记、配对繁殖计划、不同种群之间的个体交换，全是在和奠基者效应赛跑。

\section{基因流：群体之间的桥}

前面讲的都是单个群体内部的故事。可群体之间并不是孤岛——只要有个体迁移、并且在新家留下后代，等位基因就跟着搬了家。这叫\textbf{基因流}。

基因流的“搬运工”五花八门。候鸟把等位基因带上几千公里；枫树、玉米的花粉乘风飘过几块田；洄游的鲑鱼偶尔会“走错家门”游进相邻的河流；人类就更不必说——历史上每一次迁徙、通商和远航，都是一次大规模基因流，今天几乎找不到与世隔绝的人类群体。

基因流的进化效果可以概括成一句话：\textbf{它让不同群体变得更像彼此。}两个池塘的鱼，哪怕环境不同、选择方向不同，只要一直有少数个体来往，它们的基因库就很难真正分道扬镳。反过来，一旦基因流被切断（山升起、海淹没、河流改道），两个群体就开始各自漂变、各自适应，差异越积越多——这正是下一章“物种怎样分开”的起点，先把这颗扣子留在这里。

基因流也能被人类用来“救急”，这是它在保护生物学里最漂亮的应用。1990 年代，美国佛罗里达美洲狮只剩二三十只，长期近交让它们出现心脏缺陷、精子畸形和尾巴打结。1995 年，保护人员做了一次大胆决定：从得克萨斯引进 8 只美洲狮。这不是“放老虎归山”式的莽撞——两个种群本就是同一亚种，历史上也曾连通。结果：混血后代的缺陷明显减少，种群数量回升。8 只外来个体带来的基因流，把一潭正在变质的死水重新接回了河流。当然，基因流也有另一面：外来个体可能冲淡本地群体好不容易积累的适应（这叫“远交衰退”），所以每一次人为的基因救援都要先掂量两边群体分开了多久、差异有多大。

\section{性选择：活得久，不如留得下}

达尔文在《物种起源》出版十二年后，又专门写了一本书处理一个让他头疼的问题：孔雀的尾巴。那身华丽累赘的羽毛，怎么看都不利于生存——飞不快、藏不住，简直是给狐狸递请柬。如果“适者生存”是唯一规则，这种尾巴早该被淘汰干净。

答案在于：\textbf{自然选择筛选的是“留下多少后代”，而不是“活得多体面”。}一只雄孔雀哪怕寿命短一点，只要尾巴足够吸引雌孔雀，留下的后代就比朴素同伴多——华丽的等位基因照样扩散。雌孔雀的“偏好”本身也有遗传基础，于是偏好和尾巴在一代代中互相加码，越滚越夸张。这套机制叫\textbf{性选择}。

严格地说，很多生物学家把性选择看作自然选择的一个组成部分（毕竟都是繁殖成功的差异）。但把它单独拎出来讲有现实意义：它提醒我们，\textbf{“被选择”不等于“更健康、更强壮、更适应环境”}。雄鹿头顶几十斤重的角、园丁鸟费时几周搭建的求偶亭、深海鮟鱇鱼雄鱼干脆“嫁”进雌鱼身体——这些性状在生存账本上都是负资产，在繁殖账本上却是暴利。进化算的是总账。

性选择还能反过来影响前面讲的机制：少数“赢家”雄性占有大量交配机会时，实际参与繁殖的个体数远小于群体总数，等于人为地把群体变小——漂变的音量又被调大了。进化的各台发动机，从来不是各自独立运转的。

\section{中性理论：分子层面，大多数变化“无所谓”}

到这里，一个奇怪的问题浮出水面。20 世纪 60 年代，科学家第一次能直接比较不同物种的蛋白质序列，结果发现：分子层面积累的差异，多得远超“适者生存”能解释的量。如果每一个被保留下来的差异都是自然选择的功劳，那选择的“工作量”也太大了——一个群体根本负担不起那么多“优胜劣汰”。

日本群体遗传学家木村资生在 1968 年给出了一个当时颇为离经叛道的回答：\textbf{分子层面的绝大多数突变，既不利也不害，是中性的。}它们的命运不取决于选择，而取决于漂变。想想第 72 章的知识就明白这很合理：遗传密码有简并性，很多碱基替换根本不改变氨基酸（沉默突变）；基因组里还有大片不编码蛋白质的序列，在里面改几个字母，蛋白质机器毫无察觉。这些中性变异在群体里随漂变起起落落，偶尔撞上大运被固定下来——分子进化的大部分脚印，其实是这么踩出来的。

注意，中性理论不是说“自然选择不重要”，更不是说“性状都是随机的”。翅膀、眼睛、解毒酶这些明显影响生存繁殖的特征，依然牢牢握在选择手里。木村说的是：把镜头推进到 DNA 序列的颗粒度，你会发现大量变化发生在选择的“雷达盲区”里。选择决定了一部小说的情节主线，中性漂变决定了字里行间的无数标点。

\begin{zhengju}
中性理论是怎样被接受又不断被修正的？这个故事值得细讲，因为它几乎是一份“科学如何运转”的标准样本。

先是提出问题。1960 年代初，祖克坎德尔和鲍林比较了不同动物的血红蛋白序列，发现一个规律：两个物种分开得越久，序列差异越多，而且累积速度大致均匀，像一座“滴答作响的分子钟”。为什么进化会这么守时？如果每次差异固定都是自然选择推动的，速度应该随环境变化忽快忽慢，不该这么均匀。

再是给出解释和替代解释。木村资生 1968 年提出：若大多数固定下来的突变是中性的，固定速率就直接等于突变率——而突变率相对稳定，钟自然就走得匀。这是一个干净、可以检验的预言。反对意见也随之而来：也许大多数突变其实是轻微有害的，只是还没来得及被清除？也许蛋白质的序列变化大多经历过选择，只是我们看不出选择的理由？争论持续了十几年，双方都在比对新测的序列、改进统计方法。

最后是修正与共存。越来越多的基因组数据表明：中性理论在解释非编码区和沉默位点的大尺度模式时非常出色，但许多蛋白质里确实存在被选择反复“清扫”或“保留”的痕迹。今天的教科书把两者并排摆放：\textbf{中性漂变是背景噪音，自然选择是穿过噪音的信号}——而“分子钟”也相应地被打了折扣：不同基因、不同类群的钟走得快慢不一，用之前必须先校准。木村晚年说，他的理论不是取代达尔文，而是给达尔文补上了分子这一层。科学就是这样：没有谁是最后的赢家，只有越来越清楚的边界。
\end{zhengju}

分子钟虽然需要校准，但已经足够好用。它的逻辑很直白：如果中性突变以大致稳定的速率积累，那么两个物种 DNA 序列的差异量，就正比于它们从共同祖先分开的时间。数出差异，除以速率，就能估算“分家”的年代。用它核对化石记录，多数时候对得上；对不上的地方，往往藏着有趣的故事（比如某段DNA受到了选择，或两类群的突变率不同）。第 78 章讲系统发育树时，我们还会用到这套“数差异”的手艺。

\section{一张表，看清这几台发动机}

把这一章的四股力量并排摆开（突变是原料供应，单独说过了）：

\begin{table}[htbp]
\centering
\begin{tabular}{llll}
\toprule
机制 & 改变频率的方式 & 小群体中 & 会产生适应吗 \\
\midrule
自然选择 & 差异繁殖成功，有方向 & 依然有效，但易被噪音盖过 & 会，适应的唯一来源 \\
遗传漂变 & 随机抽样误差，无方向 & 威力巨大 & 不会，但可固定任何变异 \\
基因流 & 群体间搬运等位基因 & 取决于迁移量 & 间接：输入新变异 \\
性选择 & 交配机会不平等 & 赢家通吃，缩小有效群体 & 产生“繁殖适应”，未必利生存 \\
\bottomrule
\end{tabular}
\end{table}

注意最后一列：在这张表里，\textbf{能产生适应的只有选择}（包括性选择那种别扭的适应）。漂变和基因流只负责改变频率，不管方向是好是坏。这是本章最重要的一行字：进化不等于适应，进化比适应大得多。

下面这张图是一次豆子模拟可能得到的曲线——某个等位基因的频率在 30 代里纯粹靠运气游走，最后撞上 0 或 100\% 中的一端就再也回不来。这是人为的表示法：真实群体没有坐标轴，但频率的起落就是这样没有剧本。

\begin{center}
\begin{tikzpicture}[scale=0.9]
\draw[->] (0,0) -- (11,0) node[right] {世代};
\draw[->] (0,0) -- (0,4.4) node[above] {频率};
\draw[dashed] (0,2) -- (10.5,2) node[right] {50\%};
\draw (0,0) node[left] {0} (0,4) node[left] {100\%};
\draw[thick] (0,2)--(0.7,2.3)--(1.4,1.9)--(2.1,2.2)--(2.8,1.6)--(3.5,1.8)--(4.2,1.3)--(4.9,1.5)--(5.6,1.0)--(6.3,1.2)--(7.0,0.7)--(7.7,0.9)--(8.4,0.4)--(9.1,0.2)--(9.8,0);
\draw (9.8,0) node[below right] {丢失};
\end{tikzpicture}
\end{center}

\section{这套解释什么时候会失手}

本章的模型有清晰的适用边界，也诚实地告诉你三件事。

第一，\textbf{真实群体几乎不满足任何一个理想假设}。豆子模拟里的群体大小恒定、随机交配、没有迁移、没有选择——自然界里这四个条件能同时满足的几乎没有。模型的价值不在于“像”，而在于给出基准：先把所有机制关掉，看运气本身能干出多大的事；实测数据偏离这个基准多少，就是其他机制存在的证据。

第二，\textbf{漂变和选择在真实数据里常常缠在一起，很难拆开}。一个等位基因频率上升，既可能是它带来了优势，也可能只是运气好，还可能是它恰好和某个真正有利的基因“搭了便车”（连锁在同一条染色体上）。判断谁是谁，需要大样本、多基因位点和统计模型——这是现代群体遗传学每天都在啃的硬骨头，没有速成的判别法。

第三，\textbf{分子钟不是节拍器}。不同基因的钟快慢不同（受功能约束强弱影响），不同类群的钟也不同（世代时间短的物种突变积累快），甚至同一基因在不同年代也会变速。用分子钟推算年代，必须多基因交叉验证、用化石校准，单拿一个基因就算出“精确分家时间”的说法，都值得警惕。

\begin{wuqu}
“既然存在进化，生物身上每个特征都一定有它的用处。”——这叫“适应主义”的过度延伸，是流传极广的误解。反例就在本章里：平格拉普岛民的全色盲没有任何用处，是台风筛选出的运气；你身上 ABO 血型在世界各地人群中的频率差异，大部分也能用漂变和人群历史解释，而非各地环境“需要”不同的血型。还有一些特征是其他适应的副产品或历史遗留，例如男人不长乳腺却有乳头——因为胚胎发育的同一套图纸要造两种性别，专门删掉乳头的代价可能高于留着它。遇到任何性状，正确的问题不是“它有什么用”，而是“它可能是怎么来的”：选择、漂变、基因流、历史偶然，都是候选答案。
\end{wuqu}

\section{回到那罐弹珠}

现在兑现开场的承诺。一罐 50 红 50 黑的弹珠，每代只留 10 颗繁殖——结局几乎注定：经过足够多代，罐子会变成全红或全黑，而且一旦如此就永远定型（除非有突变补充新颜色）。这不是因为某种颜色“更好”，而是每一次抽样都在累积微小的运气，而运气的终点只有两个字：\textbf{固定}。

哪个颜色赢？无法预测。会漂多久？群体越小越快。这罐弹珠就是过去几亿年里无数小岛物种、灾后残存种群的缩影：它们身上许多特征，不是被环境精心挑选的勋章，而是被运气随手贴上的标签。想读懂一个物种，只问“环境选了什么”是不够的，还要问“它经历过几次瓶颈、祖先从哪迁来、和谁通过婚、群体曾经多小”。

\section{分歧之后，会发生什么}

这一章我们一直在一个群体内部看频率的涨落。但请把镜头拉远：如果两个群体之间的桥断了——基因流归零——它们就会各自漂变、各自积累突变、各自经受各自环境的选择。差异一代代叠加，总有一天多到一个临界点：就算重新相遇，也无法再生育后代。到那一刻，一棵树分出了两根不再相交的枝条。

这就是下一章的主题：物种怎样分开。我们会遇到比“能不能交配”更微妙的问题——环物种、杂交带、根本不走有性生殖的细菌——还会用到本章磨好的工具：中性理论提供的分子钟，和漂变告诉我们的“分歧不完全等于适应”。而再往后第 79 章，我们会潜回基因组内部，看进化这位工匠如何用复制、借用和改造旧零件，拼凑出生命的新花样。

\begin{tiaozhan}
漂变有一个优美得近乎反直觉的定量结论，供有兴趣的你推导。在一个理想群体（$N$ 个个体，$2N$ 份基因拷贝）里，一个当前频率为 $p$ 的中性等位基因，最终固定下来的概率恰好就是 $p$ 本身——稀有变体也可能赢，只是概率小。由此还能推出：一个刚产生的新中性突变（在群体里只有 1 份拷贝，频率 $1/(2N)$），固定概率就是 $1/(2N)$。而每一代新产生的中性突变总数，是整个群体的拷贝数乘以每个拷贝的突变率 $u$，即 $2Nu$ 个。两者相乘：每代被固定的中性突变数 $= 2Nu \times \frac{1}{2N} = u$。\textbf{群体大小 $N$ 竟然被约掉了}——中性突变的固定速率就等于突变率，与群体规模无关。这正是分子钟走得大致均匀的数学根基，也是木村资生当年最有说服力的那张底牌。推导只用了一个假设：每个拷贝传给下一代的机会均等。跳过去完全不影响主线，但值得看一眼——一个含 $N$ 的量最终被约掉，是物理学和生物学里最美的那类惊喜。
\end{tiaozhan}

\section*{练一练}

\begin{sikaoti}
\item 画一条信息流图：从“DNA 复制出错”开始，经过“新等位基因进入群体”，分岔出三条路径（被选择保留、被漂变固定、丢失），每条路径标注依赖的条件（环境、群体大小、运气）。哪条路径可以脱离环境发生？
\item 一个 20 只个体的鸟群和一个 20000 只个体的鸟群，各有一个频率 50\% 的中性等位基因。用本章的公式估算两个群体下一代频率的标准差，并解释为什么小鸟群更可能在几代之内丢失这个等位基因。
\item 平格拉普岛的全色盲和佛罗里达美洲狮的基因救援，一个是奠基者效应的后果，一个是基因流的应用。用“弹珠罐”的语言各画一幅示意图，说明两个故事中等位基因频率分别因什么而改变。图旁注明：弹珠只是表示法，真实基因库远比两种颜色复杂。
\item 有人说：“孔雀尾巴是进化的证据，说明生物会变得越来越完美。”请用本章和上一章的概念，写三句话反驳，要求至少用到“适应只针对当时环境”和“性选择”两个要点。
\item 木村资生说分子层面的大多数进化变化是中性的。这和“长颈鹿的脖子是选择出来的”矛盾吗？设计一个思想实验（比较编码区与沉默位点的序列差异积累速度），说明你如何用数据同时检验这两句话。
\item 朱鹮从 7 只恢复到数千只，为什么保育工作者仍然担心它的未来？从等位基因多样性的角度回答，并给出一个你认为可行的保护措施。
\end{sikaoti}
